% chktex-file 12
% chktex-file 36
\section{Identification with the \texorpdfstring{$2$}{2}-modular basic spin representation}\label{sec:spin}

This section identifies $W_n$ in two ways. Section~\ref{subsec:specht} proves the identification in the title, $W_n\cong D^{(n+1,n-1)}$, from James's definitions of Specht modules, and it needs no Clifford algebra in characteristic $0$. Sections~\ref{subsec:cover} to~\ref{subsec:reduction} work with the Clifford algebra over an algebraically closed field of characteristic $0$ through which the double cover $\St_{2n}$ acts, and prove that the reduction modulo $2$ of each of its finite-dimensional simple modules is $k\otimes_{\F}W_n$, for every lattice and every reduction map.

\subsection{The double cover and the root form}\label{subsec:cover}

For $m\ge2$ let $\St_m$ be the group generated by $z,t_0,\dots,t_{m-2}$ with the relations
\[
z^2=1,\quad z\text{ central},\quad t_j^2=z,\quad {(t_{j}t_{j+1})}^3=z,\quad {(t_{i}t_j)}^2=z\ (j\ge i+2).
\]
The map $t_j\mapsto\tr j{j+1}$, $z\mapsto1$ is a surjection $\pi\colon\St_m\to S_m$. This is the presentation used by Maas~\cite{Maas10}, with his generators $t_1,\dots,t_{m-1}$ numbered from $0$. The ordinary basic spin representations were constructed by Schur~\cite{Schur11}, as Maas recalls in his introduction~\cite{Maas10}.

Let $\mathcal O$ be a commutative ring. On $\mathcal O^{m-1}$ with standard basis $b_0,\dots,b_{m-2}$ put
\[
Q_{\mathrm S}(x)=-\sum_{j=0}^{m-2}x_j^2+\sum_{j=0}^{m-3}x_{j}x_{j+1},
\]
the negative of the root form of type $A_{m-1}$. Then $Q_{\mathrm S}(b_j)=-1$, and the polar form $Q_{\mathrm S}(x+y)-Q_{\mathrm S}(x)-Q_{\mathrm S}(y)$ takes the values $-2$ on the diagonal, $1$ on adjacent basis vectors and $0$ otherwise. In the Clifford algebra $\Cl_{\mathcal O}(Q_{\mathrm S})$ the elements $T_j=\iota(b_j)$ satisfy $T_j^2=-1$, $T_{j}T_{j+1}+T_{j+1}T_j=1$ and $T_{i}T_j+T_{j}T_i=0$ for $|i-j|\ge2$.

\begin{lemma}\label{lem:schurhom}
Let $\mathcal A$ be a ring and $T_0,\dots,T_{m-2}\in\mathcal A$ with $T_j^2=-1$, adjacent anticommutators $1$ and distant anticommutators $0$. Then $t_j\mapsto T_j$, $z\mapsto-1$ defines a homomorphism $\St_m\to\mathcal A^\times$.
\end{lemma}

\begin{proof}
The element $-1$ is central with square $1$, and $T_j^2=-1$. Let $a^2=b^2=-1$ and $c=ab$. If $ab+ba=1$, then $c^2=a(ba)b=a(1-ab)b=c-a^{2}b^2=c-1$, so $c^3=c^2-c=-1$. If $ab+ba=0$, then $c^2=-a^{2}b^2=-1$. So the images satisfy the defining relations.
\end{proof}

Taking $T_j=\iota(b_j)$ in $\Cl_{\mathcal O}(Q_{\mathrm S})$ gives a homomorphism $\St_m\to{\Cl_{\mathcal O}(Q_{\mathrm S})}^\times$, which we call the \emph{Clifford homomorphism}. In characteristic $2$, for $m=2n$, Lemma~\ref{lem:schurhom} applies to the elements $\iota(v_{j,j+1})$, $0\le j\le2n-2$, of $\Cl(U_n,\qf)$, since ${\iota(v_{j,j+1})}^2=\qf(v_{j,j+1})=1=-1$ and the anticommutator of $\iota(v_{i,i+1})$ and $\iota(v_{j,j+1})$ is $\bq(v_{i,i+1},v_{j,j+1})$, which is $1$ for $|i-j|=1$ and $0$ for $|i-j|\ge2$. It gives a homomorphism $\St_{2n}\to{\Cl(U_n,\qf)}^\times$ with $t_j\mapsto\iota(v_{j,j+1})$ and $z\mapsto1$, and composed with the Clifford action on $W_n$ it equals $\rho_W\circ\pi$, because $\iota(v_{j,j+1})$ acts as $\Gamma_{j,j+1}=\rho_W(\tr j{j+1})$.

\subsection{An upper bound in characteristic \texorpdfstring{$0$}{0} and a lower bound in characteristic \texorpdfstring{$2$}{2}}\label{subsec:bounds}

Proposition~\ref{prop:jsystem} and Theorem~\ref{thm:char0} bound the dimension of a Clifford spin module from above. Theorem~\ref{thm:vacuum} shows that every nonzero $\Cl(U_n,\qf)$-module in characteristic~$2$ contains a copy of $k\otimes_{\F}W_n$.

\begin{proposition}\label{prop:jsystem}
Let $K$ be a field with $2\ne0$ that contains an element $I$ with $I^2=-1$. Let $Y\ne0$ be a finite-dimensional $K$-vector space and $J_0,\dots,J_{2p}$ pairwise anticommuting operators on $Y$ with $J_r^2=1$. If the only subspaces of $Y$ stable under every $J_r$ are $0$ and $Y$, then $\dim_{K}Y\le2^p$.
\end{proposition}

\begin{proof}
For $i<p$ put $J_i^+=J_{2i}+IJ_{2i+1}$ and $J_i^-=J_{2i}-IJ_{2i+1}$. Then
\[
{(J_i^+)}^2={(J_i^-)}^2=0,\qquad J_i^+J_i^-+J_i^-J_i^+=4,
\]
and for $i\ne j$ the operators $J_i^\pm$ anticommute with $J_j^\pm$ and with $J_{2p}$.

Each $J_i^+$ anticommutes with the others, so it preserves their joint kernel. Since ${(J_i^+)}^2=0$, it has a nonzero kernel on every nonzero subspace that it preserves. Induction on $i$ gives a nonzero joint kernel of $J_0^+,\dots,J_{p-1}^+$. The operator $J_{2p}$ preserves this joint kernel and $J_{2p}^2=1$, so the joint kernel contains a vector $y_0\ne0$ with $J_{2p}y_0=\pm y_0$.

Let $A_0=Ky_0$ and $A_{j+1}=A_j+J_j^-A_j$. Then $A_p$ is spanned by the $2^p$ vectors $J_{i_1}^-\cdots J_{i_r}^-y_0$ with $i_1<\cdots<i_r$, so $\dim A_p\le2^p$.

Moving $J_i^+$, $J_i^-$ or $J_{2p}$ to the right through such a product, using the anticommutation relations, ${(J_i^-)}^2=0$, $J_i^+J_i^-=4-J_i^-J_i^+$, $J_i^+y_0=0$ and $J_{2p}y_0=\pm y_0$, shows that $A_p$ is stable under all of them. Since $J_{2i}=(J_i^++J_i^-)/2$ and $J_{2i+1}=(J_i^+-J_i^-)/(2I)$, $A_p$ is stable under every $J_r$, so $A_p=Y$.
\end{proof}

\begin{definition}\label{def:basicspin}
Let $K$ be an algebraically closed field of characteristic $0$ and $n\ge1$. A \emph{Clifford spin module} of $\St_{2n}$ over $K$ is a finite-dimensional simple module $Y$ over $\Cl_K(Q_{\mathrm S})$, with $Q_{\mathrm S}$ on $K^{2n-1}$, regarded as a representation of $\St_{2n}$ through the Clifford homomorphism.
\end{definition}

\todo[inline]{Source needed. A read reference, with locator, stating that the Clifford spin modules of Definition~\ref{def:basicspin} are the two associate basic spin representations of $\St_{2n}$ of Schur, as used by Maas~\cite[Thm.~1]{Maas10}. Candidates are Hoffman and Humphreys (1992) and Stembridge (1989), neither verified. Until then the results below are stated for Clifford spin modules.}

\begin{theorem}\label{thm:char0}
Let $K$ be algebraically closed of characteristic $0$ and $n\ge1$. Then the polar form of $Q_{\mathrm S}$ on $K^{2n-1}$ is nondegenerate, and every Clifford spin module $Y$ of $\St_{2n}$ satisfies $\dim_{K}Y\le2^{n-1}$.
\end{theorem}

\begin{proof}
Suppose the polar form vanishes on $x$ against every $b_t$. This reads $-2x_t+x_{t-1}+x_{t+1}=0$ for all $t$, with $x_{-1}=x_{2n-1}=0$. The equations for $t=0,\dots,2n-3$ give $x_t=(t+1)x_0$ by induction, and the equation for $t=2n-2$ then reads $-2(2n-1)x_0+(2n-2)x_0=-2n\,x_0=0$. So $x=0$ in characteristic $0$. Over an algebraically closed field of characteristic $0$ a nondegenerate quadratic form has a basis $w_0,\dots,w_{2n-2}$ that is orthogonal for the polar form and has $Q_{\mathrm S}(w_r)=1$. The operators $J_r=\iota(w_r)$ on $Y$ are pairwise anticommuting involutions that generate the action of $\Cl_K(Q_{\mathrm S})$, so Proposition~\ref{prop:jsystem} with $p=n-1$ applies.
\end{proof}

\begin{theorem}\label{thm:vacuum}
Let $k$ be a field of characteristic $2$, $n\ge1$, and let $M\ne0$ be a $k$-vector space with a ring homomorphism $\alpha\colon\Cl(U_n,\qf)\to\End_k(M)$.
\begin{enumerate}
\item There is $m_0\ne0$ with $\alpha(\iota(v_{2a,2a+1})+1)m_0=0$ for every block $a$.
\item There is an injective $k$-linear map $\eta\colon k\otimes_{\F}W_n\to M$ that intertwines the Clifford action on $k\otimes_{\F}W_n$ with $\alpha$.
\item If $\dim_{k}M=2^{n-1}$, then $\eta$ is an isomorphism.
\end{enumerate}
\end{theorem}

\begin{proof}
The elements $N_a=\iota(v_{2a,2a+1})+1$ satisfy $N_a^2=\qf(v_{2a,2a+1})+1=0$, and they commute because distinct blocks give orthogonal vectors. A nonzero space on which $N_0,\dots,N_{a-1}$ vanish is stable under $N_a$ and meets $\ker N_a$, since $N_a^2=0$. This gives (i). For (ii) put $u_S=\alpha\bigl(\prod_{g\in S}\iota(v_{2g+1,2g+2})\bigr)m_0$ for $S\subseteq G$. The relations used in Proposition~\ref{prop:Bw} hold in $\Cl(U_n,\qf)$, so $\alpha(\iota(v_{2a,2a+1}))$ and $\alpha(\iota(v_{2g+1,2g+2}))$ act on the $u_S$ by the same formulas as $B_a$ and the gap transpositions act on the $w_S$. Hence the $k$-linear map $\eta$ with $\eta(w_S)=u_S$ intertwines the generators of $\Cl(U_n,\qf)$, and therefore the whole algebra. The operators $\partial_g$ of Proposition~\ref{prop:lower}, formed inside $\Cl(U_n,\qf)$, satisfy $\partial_{g}u_S=[g\in S]u_{S\setminus\{g\}}$. Suppose $\sum c_{S}u_S=0$ with some $c_S\ne0$, and choose $S_0$ maximal under inclusion with $c_{S_0}\ne0$. Applying $\prod_{g\in S_0}\partial_g$ leaves $c_{S_0}m_0=0$, a contradiction. So the $u_S$ are independent. Part (iii) follows from (ii) by dimension.
\end{proof}

\subsection{Lattices and reduction modulo \texorpdfstring{$2$}{2}}\label{subsec:reduction}

Let $K$ be algebraically closed of characteristic $0$ and $Y$ a Clifford spin module of $\St_{2n}$. Fix a basis of $Y$, indexed by a finite set $\Omega$, in which the matrix $[T_j]$ of $T_j=\iota(b_j)$ has entries in a subring $\mathcal O\subseteq K$ for $0\le j\le2n-2$. We call the $\mathcal O$-span of such a basis a \emph{lattice} in $Y$. Let $k$ be a field of characteristic $2$ and $\psi\colon\mathcal O\to k$ a ring homomorphism, and let $\overline{T}_j$ be the matrix obtained by applying $\psi$ to the entries of $[T_j]$. Since $\psi$ is a ring homomorphism, $\overline{T}_j^2=\psi(-1)=1$, $\overline{T}_j\overline{T}_{j+1}+\overline{T}_{j+1}\overline{T}_j=1$, and the distant anticommutators vanish. By Lemma~\ref{lem:reducehom}, $\iota(v_{j,j+1})\mapsto\overline{T}_j$ makes $\overline{Y}=k^\Omega$ a $\Cl(U_n,\qf)$-module, and by Lemma~\ref{lem:schurhom} it is a representation of $\St_{2n}$ in which $t_j$ acts as $\overline{T}_j$. This is the reduction of the matrices of $\St_{2n}$ on $Y$.

\begin{theorem}\label{thm:reduction}
Let $n\ge1$, and let $Y$, $\Omega$, $\mathcal O$, $\psi$ and $\overline{Y}$ be as above.
\begin{enumerate}
\item $|\Omega|=\dim_{K}Y=2^{n-1}$.
\item There is a $k$-linear isomorphism $k\otimes_{\F}W_n\to\overline{Y}$ that intertwines the two $\Cl(U_n,\qf)$-actions and the two $\St_{2n}$-actions, where $\St_{2n}$ acts on $k\otimes_{\F}W_n$ through $\pi\colon\St_{2n}\to S_{2n}$.
\item $\overline{Y}$ is irreducible under the matrices $\overline{T}_j$.
\end{enumerate}
\end{theorem}

\begin{proof}
By Theorem~\ref{thm:vacuum}(ii), $k\otimes_{\F}W_n$ embeds in $\overline{Y}$, so $2^{n-1}\le|\Omega|=\dim_{K}Y$. Theorem~\ref{thm:char0} gives $\dim_{K}Y\le2^{n-1}$. This proves (i), and Theorem~\ref{thm:vacuum}(iii) gives the $\Cl(U_n,\qf)$-isomorphism. The $\St_{2n}$-actions on both sides factor through the homomorphism $\St_{2n}\to{\Cl(U_n,\qf)}^\times$ of Section~\ref{subsec:cover}, and on $k\otimes_{\F}W_n$ that homomorphism followed by the Clifford action is $\rho_k\circ\pi$. For (iii), a subspace of $\overline{Y}$ stable under every $\overline{T}_j$ is stable under all of $\Cl(U_n,\qf)$, because the $\iota(v_{j,j+1})$ generate it. Its preimage in $k\otimes_{\F}W_n$ is then stable under every $\rho_k(\sigma)$, and Theorem~\ref{thm:burnside} applies.
\end{proof}

\begin{proposition}\label{prop:lattice}
Let $K$ be algebraically closed of characteristic $0$.
\begin{enumerate}
\item There is a valuation subring $\mathcal V\subseteq K$ with $2$ in its maximal ideal.
\item Every finite-dimensional $\Cl_K(Q_{\mathrm S})$-module has a basis in which every $T_j$ has entries in $\mathcal V$.
\item A Clifford spin module of $\St_{2n}$ over $K$ exists.
\end{enumerate}
\end{proposition}

\begin{proof}
(i) The ideal $2\Z$ of the subring $\Z\subseteq K$ is proper, and every proper ideal of a subring of $K$ lies in the nonunits of some valuation subring of $K$. (ii) Let $y_1,\dots,y_d$ be a basis of the module and let $L$ be the $\mathcal V$-span of the vectors $T_{j_1}\cdots T_{j_r}y_l$ with $j_1<\cdots<j_r$ and $1\le l\le d$. Since $T_j^2=-1$ and $T_{i}T_j+T_{j}T_i\in\Z$, left multiplication by any $T_t$ maps an ordered monomial into the $\Z$-span of ordered monomials. So $L$ is stable under every $T_t$. It is finitely generated and torsion-free over the valuation ring $\mathcal V$, hence free, and it spans the module over $K$, so a $\mathcal V$-basis of $L$ is a $K$-basis. (iii) The algebra $\Cl_K(Q_{\mathrm S})$ is spanned by ordered monomials, so it is finite-dimensional. It is nonzero, and a minimal nonzero left ideal is a finite-dimensional simple module.
\end{proof}

\todo[inline]{Sources needed in the proof of Proposition~\ref{prop:lattice}. In (i), a citation with locator for Chevalley's extension theorem, since Bourbaki, Commutative Algebra Ch.~VI, is not verified. In (ii), a citation with locator for the freeness of a finitely generated torsion-free module over a valuation ring. In (iii), a source for $\Cl_K(Q_{\mathrm S})\ne0$, such as the ordered-monomial basis of a Clifford algebra, or an explicit nonzero module.}
\medskip

By Proposition~\ref{prop:lattice}, every Clifford spin module of $\St_{2n}$ has a basis as in Theorem~\ref{thm:reduction}, with $\mathcal O=\mathcal V$ and $\psi$ the map to the residue field $k=\mathcal V/\mathfrak m_{\mathcal V}$, which has characteristic $2$. Theorem~\ref{thm:reduction} says that for every such lattice and every such $\psi$ the reduction is $k\otimes_{\F}W_n$, which is irreducible, and that every Clifford spin module $Y$ has $\dim_{K}Y=2^{n-1}$. The results of Wales and Benson recalled in Section~\ref{subsec:prior} concern Schur's basic spin representations, and their comparison with Theorem~\ref{thm:reduction} rests on the identification noted after Definition~\ref{def:basicspin}.

\subsection{Polytabloids, the first Jantzen layer of \texorpdfstring{$S^{(n,n)}$}{S\textasciicircum(n,n)} and \texorpdfstring{$D^{(n+1,n-1)}$}{D\textasciicircum(n+1,n-1)}}\label{subsec:specht}

The proofs of this subsection use the definitions of James's lecture notes~\cite{James78}, the results of James quoted in Section~\ref{subsubsec:james}, and Sections~\ref{sec:matchings} to~\ref{sec:presentations} for $W_n$. Within this subsection, the isomorphism $W_n\cong D^{(n+1,n-1)}$ of Section~\ref{subsubsec:nearmatch} (Lemma~\ref{lem:nearmatch} and Theorem~\ref{thm:step4}, apart from its last isomorphism) uses only Section~\ref{subsubsec:james}. Sections~\ref{subsubsec:fourswitch} to~\ref{subsubsec:jantzen} identify $W_n$ with the first Jantzen layer of $S^{(n,n)}$. Throughout, $V=\{0,1,\dots,2n-1\}$ with its natural order, and $\M_n=\M(V)$.

\subsubsection{James's definitions}\label{subsubsec:james}

Let $\mu$ be a partition of $m$ and let $\Bbbk$ be a field, $\Z$ or $\Z_{(2)}$. A $\mu$-tableau is a filling of the Young diagram of $\mu$ by the elements of an $m$-element set. James uses $\{1,\dots,m\}$, and we use $V$ below, which changes nothing in what follows. The tabloid $\{t\}$ of a tableau $t$ is its class under permutations of the entries within rows. The permutation module $M^\mu_\Bbbk$ is the free $\Bbbk$-module on the $\mu$-tabloids~\cite[4.1]{James78}. James writes permutations on the right. We write them on the left, so that $\sigma\in S_V$ sends $t$ to the tableau $\sigma t$ obtained by applying $\sigma$ to every entry, and $\sigma\{t\}=\{\sigma t\}$. Let $C_t$ be the column stabilizer of $t$. In this notation the polytabloid of $t$ is
\[
e_t=\sum_{\gamma\in C_t}\mathrm{sgn}(\gamma)\,\gamma\{t\},
\]
and the Specht module $S^\mu_\Bbbk$ is the $\Bbbk$-span of the polytabloids in $M^\mu_\Bbbk$~\cite[4.3]{James78}. Since $C_{\sigma t}=\sigma C_t\sigma^{-1}$, writing each element of $C_{\sigma t}$ as $\sigma\gamma\sigma^{-1}$ with $\gamma\in C_t$ gives $\sigma e_t=e_{\sigma t}$. The module $M^\mu_\Bbbk$ carries the symmetric $S_V$-invariant bilinear form $\pr{\,,}$ for which the tabloids are orthonormal~\cite[\S4]{James78}, and $X^\perp$ denotes the orthogonal of a subspace $X$ of $M^\mu_\Bbbk$. When $\mu$ has two rows, a tabloid is determined by the set $R$ of entries in its first row, and we write $\{R\}$.

We use the following results of James. For every field $\Bbbk$, $S^\mu_\Bbbk/(S^\mu_\Bbbk\cap S^{\mu\perp}_\Bbbk)$ is zero or absolutely irreducible, and if it is nonzero then $S^\mu_\Bbbk\cap S^{\mu\perp}_\Bbbk$ is the unique maximal submodule of $S^\mu_\Bbbk$~\cite[Thm.~4.9]{James78}. If $\Bbbk$ has characteristic $p$ and $\mu$ is $p$-regular, that is, no part of $\mu$ occurs $p$ or more times, then $D^\mu=S^\mu_\Bbbk/(S^\mu_\Bbbk\cap S^{\mu\perp}_\Bbbk)$~\cite[Def.~11.2]{James78}. Let $g_\mu$ be the greatest common divisor of the values $\pr{e_t,e_{t'}}$ over $\Z$. Then $p$ divides $g_\mu$ if and only if $\mu$ is $p$-singular~\cite[10.3--10.5]{James78}.

\subsubsection{The four-switch quotient is \texorpdfstring{$S^{(n,n)}$}{S\textasciicircum(n,n)}}\label{subsubsec:fourswitch}

For $P\in\M_n$ call $R\subseteq V$ a \emph{transversal} of $P$ if $R$ contains exactly one endpoint of each edge of $P$. Put
\[
v_P=\sum_{R\text{ transversal of }P}{(-1)}^{\beta_P(R)}\{R\}\in M^{(n,n)}_\Z,
\]
where $\beta_P(R)$ is the number of edges $\{a<b\}$ of $P$ with $b\in R$.

\begin{lemma}\label{lem:vP}
Let $t_P$ be an $(n,n)$-tableau whose columns are the edges of $P$, each with its smaller entry in the first row. Then $e_{t_P}=v_P$, and every polytabloid of shape $(n,n)$ equals $\pm v_P$ for the matching $P$ formed by its columns. In particular $S^{(n,n)}_\Bbbk$ is the $\Bbbk$-span of the $v_P$, and $\sigma v_P=\pm v_{\sigma P}$ for $\sigma\in S_{2n}$.
\end{lemma}

\begin{proof}
The column stabilizer of $t_P$ is the product of the groups $\{1,\tr ab\}$ over the edges $\{a<b\}$. The first row of $\{t_P\}$ is the set of smaller endpoints. Applying the transpositions of a subset $D$ of edges gives the tabloid whose first row contains the larger endpoints of the edges in $D$ and the smaller endpoints of the others, with sign ${(-1)}^{|D|}$. This is $v_P$. Reordering the columns of a tableau changes neither $\{t\}$ nor $C_t$. Exchanging the two entries of a column of $t$ replaces $t$ by $ct$ for a transposition $c\in C_t$. Then $C_{ct}=C_t$ and $\{ct\}=c\{t\}$, so $e_{ct}=\sum_{\gamma\in C_t}\mathrm{sgn}(\gamma)\,\gamma c\{t\}$. Substituting $\gamma'=\gamma c$ turns this into $\sum_{\gamma'\in C_t}\mathrm{sgn}(\gamma'c)\,\gamma'\{t\}=\mathrm{sgn}(c)\,e_t=-e_t$. The last statement follows from $\sigma e_t=e_{\sigma t}$.
\end{proof}

Identify $\{R\}$ with the monomial $\prod_{r\in R}x_r\prod_{s\notin R}y_s$ in $\Z[x_0,\dots,x_{2n-1},y_0,\dots,y_{2n-1}]$. Distinct tabloids give distinct monomials, so this is an injective linear map, and expanding the product shows that $v_P$ maps to the product of the Pl\"ucker coordinates $p_{ab}=x_{a}y_b-x_{b}y_a$ over the edges $\{a<b\}$ of $P$. A permutation $\sigma$ acts on monomials by $x_r\mapsto x_{\sigma r}$ and $y_s\mapsto y_{\sigma s}$. For $a<b<c<d$ the three-term Pl\"ucker identity $p_{ac}p_{bd}=p_{ab}p_{cd}+p_{ad}p_{bc}$ holds in this ring. Multiplying it by the product of $p_{uw}$ over the edges $\{u<w\}$ of a perfect matching $E$ of $V\setminus\{a,b,c,d\}$ gives a relation among the images of three vectors $v_P$, and injectivity gives
\begin{equation}\label{eq:plucker}
v_{E\cup\{ac,bd\}}=v_{E\cup\{ab,cd\}}+v_{E\cup\{ad,bc\}}\qquad(a<b<c<d).
\end{equation}

Two edges $\{u<w\}$ and $\{u'<w'\}$ with distinct endpoints \emph{cross} if exactly one of $u',w'$ lies strictly between $u$ and $w$. A matching is \emph{noncrossing} if no two of its edges cross, and its crossing number $\mathrm{cr}(P)$ is the number of crossing pairs of edges.

\begin{lemma}\label{lem:crossing}
Let $a<b<c<d$ and let $\{u,w\}$ be an edge disjoint from $\{a,b,c,d\}$. The number of edges of $\{ab,cd\}$ crossed by $\{u,w\}$, and the number of edges of $\{ad,bc\}$ crossed by it, are each at most the number of edges of $\{ac,bd\}$ crossed by it. Hence both matchings on the right of~\eqref{eq:plucker} have smaller crossing number than the one on the left.
\end{lemma}

\begin{proof}
Let $r_1,r_2,r_3$ be the numbers of endpoints of $\{u,w\}$ in the open intervals $(a,b)$, $(b,c)$, $(c,d)$. The edge $\{u,w\}$ crosses an edge $\{u'<w'\}$ when exactly one of $u,w$ lies in $(u',w')$. So it crosses $ac$ exactly when $r_1+r_2=1$, $bd$ when $r_2+r_3=1$, $ab$ when $r_1=1$, $cd$ when $r_3=1$, $ad$ when $r_1+r_2+r_3=1$, and $bc$ when $r_2=1$. The possible triples with $r_1+r_2+r_3\le2$ give the counts in Table~\ref{tab:crossing}, and in every row the last two entries are at most the first. For a pairing $X$ of $\{a,b,c,d\}$, the crossing number of $E\cup X$ is $\mathrm{cr}(E)$, plus the number of crossings between edges of $E$ and edges of $X$, plus the number of crossings inside $X$. The middle term does not increase when $\{ac,bd\}$ is replaced by either other pairing, and the last term is $1$ for $\{ac,bd\}$ and $0$ for $\{ab,cd\}$ and $\{ad,bc\}$.
\end{proof}

\begin{table}[ht]
\centering
\begin{tabular}{cccc}
\toprule
$(r_1,r_2,r_3)$ & crossings with $\{ac,bd\}$ & with $\{ab,cd\}$ & with $\{ad,bc\}$\\
\midrule
$(0,0,0)$, $(2,0,0)$, $(0,2,0)$, $(0,0,2)$ & $0$ & $0$ & $0$\\
$(1,0,0)$ & $1$ & $1$ & $1$\\
$(0,0,1)$ & $1$ & $1$ & $1$\\
$(0,1,0)$ & $2$ & $0$ & $2$\\
$(1,1,0)$ & $1$ & $1$ & $1$\\
$(0,1,1)$ & $1$ & $1$ & $1$\\
$(1,0,1)$ & $2$ & $2$ & $0$\\
\bottomrule
\end{tabular}
\caption{Crossings of an outside edge with the three pairings of $a<b<c<d$.}\label{tab:crossing}
\end{table}

For a noncrossing matching $A$ let $O(A)$ be the set of smaller endpoints of its edges. In a noncrossing matching the points strictly between the two ends of an edge are matched among themselves. So each larger endpoint $j$ is matched to the largest point of $O(A)$ below $j$ that is not matched to a larger endpoint below $j$, and $A$ is determined by $O(A)$.

\begin{theorem}\label{thm:step1}
Let $n\ge1$.
\begin{enumerate}
\item The vectors $v_A$, $A$ noncrossing, form a $\Z$-basis of $S^{(n,n)}_\Z$, and their reductions modulo $2$ form an $\F$-basis of $S^{(n,n)}_{\F}$.
\item The map $\nu\colon\F[\M_n]\to M^{(n,n)}_{\F}$, $e_P\mapsto v_P\bmod2$, is $S_{2n}$-equivariant, has image $S^{(n,n)}_{\F}$ and kernel $R_4$. Hence $\Delta_{2n}(0)\cong S^{(n,n)}_{\F}$ as $\F S_{2n}$-modules.
\item Let $A$ be noncrossing and $0\le i\le2n-2$. If $\{i,i+1\}\in A$, then $\tr i{i+1}v_A=-v_A$. Otherwise let $a$ and $b$ be the partners of $i$ and $i+1$ in $A$, and let $A'$ be obtained from $A$ by replacing the edges $\{i,a\}$ and $\{i+1,b\}$ with $\{i,i+1\}$ and $\{a,b\}$. Then $A'$ is noncrossing and $\tr i{i+1}v_A=v_A+v_{A'}$. Consequently, through $\nu$, the operator $e_i=1+\tr i{i+1}$ of Theorem~\ref{thm:delta} acts on the classes of the noncrossing matchings by $e_{i}A=0$ if $\{i,i+1\}\in A$ and $e_{i}A=A'$ otherwise.
\end{enumerate}
\end{theorem}

\begin{proof}
(i) Spanning. By Lemma~\ref{lem:vP}, $S^{(n,n)}_\Z$ is the $\Z$-span of the $v_P$. If $P$ has a crossing pair, it has one of the form $\{ac,bd\}$ with $a<b<c<d$, and~\eqref{eq:plucker} writes $v_P$ as a sum of two terms of smaller crossing number by Lemma~\ref{lem:crossing}. Induction on $\mathrm{cr}(P)$ writes every $v_P$ as an integer combination of noncrossing $v_A$.

Independence over every commutative ring. The set $O(A)$ is a transversal of $A$ containing no larger endpoint, so $\{O(A)\}$ has coefficient $1$ in $v_A$. Suppose $O(A)$ is a transversal of another noncrossing $B$. Each element of $O(A)$ is at least the smaller endpoint of its edge of $B$, which gives a bijection $O(A)\to O(B)$ that does not increase any element. So the sorted sequence of $O(B)$ is componentwise at most that of $O(A)$, with equality only if $O(A)=O(B)$, that is, $A=B$. Order the noncrossing matchings by a linear extension of this componentwise order. In a relation $\sum_{B}c_{B}v_B=0$, let $A$ be the first matching with $c_A\ne0$. The coefficient of $\{O(A)\}$ is $c_A$ plus terms $c_B$ with $B$ strictly earlier, so it equals $c_A\ne0$, a contradiction.

The reductions of the $v_P$ span $S^{(n,n)}_{\F}$ by Lemma~\ref{lem:vP}, each is an integer combination of noncrossing $v_A$ reduced modulo $2$, and the reductions of the noncrossing $v_A$ are independent over $\F$.

(ii) Modulo $2$, $\sigma v_P=v_{\sigma P}$, so $\nu$ is equivariant, and its image is the span of the $v_P\bmod2$, which is $S^{(n,n)}_{\F}$. Let $a,b,c,d$ be distinct and $E$ a matching of the other points. A transversal $R$ of $E\cup X$, for $X$ a pairing of $\{a,b,c,d\}$, meets $\{a,b,c,d\}$ in a $2$-set that is not an edge of $X$. Each $2$-subset of $\{a,b,c,d\}$ is an edge of exactly one of the three pairings, so it is a transversal of exactly two of them. Hence $\nu$ kills every four-switch sum and $R_4\subseteq\ker\nu$. Modulo $R_4$ the relation $e_{E\cup\{ac,bd\}}\equiv e_{E\cup\{ab,cd\}}+e_{E\cup\{ad,bc\}}$ and Lemma~\ref{lem:crossing} show that the classes of the noncrossing $e_A$ span $\F[\M_n]/R_4$. Its dimension is at most the number of noncrossing matchings, which is $\dim S^{(n,n)}_{\F}$ by (i). So the induced surjection $\F[\M_n]/R_4\to S^{(n,n)}_{\F}$ is an isomorphism.

(iii) Let $s=\tr i{i+1}$. On monomials it exchanges $x_i$ with $x_{i+1}$ and $y_i$ with $y_{i+1}$, so it fixes the Pl\"ucker coordinate of every edge of $A$ that avoids $i$ and $i+1$. If $\{i,i+1\}\in A$, then $s$ sends $p_{i,i+1}$ to $p_{i+1,i}=-p_{i,i+1}$, so $sv_A=-v_A$. Otherwise, since $A$ is noncrossing, the four points are in one of the orders $a<i<i+1<b$, $b<a<i<i+1$ or $i<i+1<b<a$. The permutation $s$ replaces the Pl\"ucker coordinates of $\{i,a\}$ and $\{i+1,b\}$ by those of $\{i+1,a\}$ and $\{i,b\}$, which cross. In each order the Pl\"ucker identity for the four sorted points writes the product of these two coordinates as the product for $\{i,a\},\{i+1,b\}$ plus the product for $\{a,b\},\{i,i+1\}$, each coordinate $p_{uw}$ taken with $u<w$. Multiplying by the coordinates of the other edges and using injectivity gives $sv_A=v_A+v_{A'}$. An edge of $A$ that crossed $\{a,b\}$ would cross $\{i,a\}$ or $\{i+1,b\}$, and no edge crosses $\{i,i+1\}$, so $A'$ is noncrossing. Reducing modulo $2$ gives the last statement.
\end{proof}

By (ii) and (iii), the classes of the noncrossing perfect matchings form a basis of $\Delta_{2n}(0)$, on which $e_i$ acts by $e_{i}A=0$ if $\{i,i+1\}\in A$ and $e_{i}A=A'$ otherwise. For $n\ge2$, since $e_i\mapsto1+\tr i{i+1}$ is an action of $\TL_{2n}(0)$ on $\Delta_{2n}(0)$ (Theorem~\ref{thm:delta}), this rule defines an action of $\TL_{2n}(0)$ on the $\F$-vector space with basis the noncrossing perfect matchings of $V$. In this paper that module is the \emph{standard module of $\TL_{2n}(0)$ with no through strands}, and $\Delta_{2n}(0)$ is this module. Drawn as diagrams with no through strands, the rule is the action of the Temperley--Lieb generators in which a closed loop gives the loop value $0$. % chktex 8

Over $\Z$ the same computation shows that $\tr i{i+1}$ acts on $S^{(n,n)}_\Z$ as $1+e_i$, where $e_{i}v_A=-2v_A$ if $\{i,i+1\}\in A$ and $e_{i}v_A=v_{A'}$ otherwise. This is the same rule at loop value $-2$, the value of Proposition~\ref{prop:jw}(ii).

\subsubsection{The Gram matrix on matchings}

\begin{lemma}\label{lem:gram}
For all $P,Q\in\M_n$, $\pr{v_P,v_Q}=\pm2^{\kappa(P\cup Q)}$.
\end{lemma}

\begin{proof}
A tabloid $\{R\}$ occurs in both $v_P$ and $v_Q$ exactly when $R$ is a transversal of both. Each component $C$ of $P\cup Q$ is an alternating cycle, a doubled edge included. A common transversal meets each edge of $C$ once, so it contains every second vertex of $C$, and there are two choices on each cycle. So there are $2^{\kappa(P\cup Q)}$ common transversals. On a cycle with $\ell$ edges of $P$ and $\ell$ edges of $Q$, the two choices differ at the chosen endpoint of each of these $2\ell$ edges. Switching the choice therefore changes $\beta_P(R)$ by $\ell$ and $\beta_Q(R)$ by $\ell$ modulo $2$, so the sign ${(-1)}^{\beta_P(R)+\beta_Q(R)}$ of the contribution of $\{R\}$ is the same for every common transversal. Hence all $2^{\kappa}$ contributions have the same sign.
\end{proof}

Every entry of the Gram matrix of $S^{(n,n)}_\Z$ is therefore even, and entries equal to $\pm2$ occur whenever $\kappa=1$. So $g_{(n,n)}=2$, consistent with James's criterion~\cite[10.3--10.5]{James78}, since $(n,n)$ is $2$-singular. In particular $S^{(n,n)}_{\F}\subseteq S^{(n,n)\perp}_{\F}$, James's quotient of $S^{(n,n)}_{\F}$ is zero, and $D^{(n,n)}$ is not defined. The module $W_n$ appears one step down the Jantzen filtration (Theorem~\ref{thm:step3}(iv)).

\subsubsection{The first Jantzen layer of \texorpdfstring{$S^{(n,n)}_{\F}$}{S\textasciicircum(n,n)}}\label{subsubsec:jantzen}

Let $\mathcal L=S^{(n,n)}_\Z\otimes\Z_{(2)}$, the $\Z_{(2)}$-span of the $v_P$ in $M^{(n,n)}_{\Z_{(2)}}$, with the restriction of $\pr{\,,}$. By Theorem~\ref{thm:step1}(i) it is free on the noncrossing $v_A$, and reduction modulo $2$ maps it onto $S^{(n,n)}_{\F}$ with kernel $2\mathcal L$. We identify $\mathcal L/2\mathcal L=S^{(n,n)}_{\F}$. The form is positive definite on $M^{(n,n)}_\Q$, so it is nondegenerate on $\mathcal L\otimes\Q$. Put
\[
\mathcal L^i=\{x\in\mathcal L:\pr{x,\mathcal L}\subseteq2^i\Z_{(2)}\},\qquad \overline{\mathcal L}^i=(\mathcal L^i+2\mathcal L)/2\mathcal L\subseteq S^{(n,n)}_{\F}.
\]
Let $\mathcal L^\#=\{y\in\mathcal L\otimes\Q:\pr{y,\mathcal L}\subseteq\Z_{(2)}\}$, which contains $\mathcal L$. Since $\mathcal L^i=\mathcal L\cap2^i\mathcal L^\#$, the $\overline{\mathcal L}^i$ form the filtration of Künzer and Nebe~\cite[Def.~1.3]{KN03} in the case $X=\mathcal L$, $Y=\mathcal L^\#$ and $\pi=2$, with the inclusion $\mathcal L\to\mathcal L^\#$ as the map. These are the terms of the Jantzen filtration of $S^{(n,n)}_{\F}$ at the prime $2$, over $\Z_{(2)}$.

\begin{theorem}\label{thm:step3}
Let $n\ge1$.
\begin{enumerate}
\item $\overline{\mathcal L}^0=\overline{\mathcal L}^1=S^{(n,n)}_{\F}$.
\item The rule $h(\bar x,\bar y)=\tfrac12\pr{x,y}\bmod2$ defines a symmetric $S_{2n}$-invariant bilinear form on $S^{(n,n)}_{\F}$, and $\overline{\mathcal L}^2=\rad h$.
\item $H_{2n}(x,y)=h(\nu x,\nu y)$ for all $x,y\in\F[\M_n]$, where $H_{2n}(x,y)=x^{\mathsf T}H_{2n}y$.
\item $\overline{\mathcal L}^1/\overline{\mathcal L}^2\cong\F[\M_n]/{\rad H_{2n}}\cong W_n$ as $\F S_{2n}$-modules, of dimension $2^{n-1}$.
\end{enumerate}
\end{theorem}

\begin{proof}
(i) By Lemma~\ref{lem:gram} every $\pr{v_P,v_Q}$ is even, so $\mathcal L^1=\mathcal L$. (ii) For $x,y\in\mathcal L$, $\pr{x,y}\in2\Z_{(2)}$ by (i). Replacing $x$ by $x+2x'$ changes $\tfrac12\pr{x,y}$ by $\pr{x',y}\in2\Z_{(2)}$, so $h$ is well defined, and it inherits symmetry and invariance. Since $2\mathcal L\subseteq\mathcal L^2$, $\bar x\in\rad h$ if and only if $\pr{x,\mathcal L}\subseteq4\Z_{(2)}$, that is, $x\in\mathcal L^2$. (iii) By Lemma~\ref{lem:gram}, $h(\nu e_P,\nu e_Q)=2^{\kappa(P\cup Q)-1}\bmod2=[\kappa(P\cup Q)=1]=H_{2n}[P,Q]$, and both sides are bilinear. (iv) Since $\nu$ is onto $S^{(n,n)}_{\F}$, (iii) gives $\rad H_{2n}=\nu^{-1}(\rad h)$, so $\nu$ induces $\F[\M_n]/{\rad H_{2n}}\cong S^{(n,n)}_{\F}/{\rad h}=\overline{\mathcal L}^1/\overline{\mathcal L}^2$, equivariantly. Theorem~\ref{thm:rank} identifies the left side with $W_n$ and gives its dimension.
\end{proof}

For $n\ge2$, Theorem~\ref{thm:presentation} gives $\rad H_{2n}=R_3$, so $\overline{\mathcal L}^2=\nu(R_3)=R_3/R_4$, which is $\pr{f_3}\Delta_{2n}(0)$ by Theorem~\ref{thm:delta}.

\subsubsection{Near-matchings and \texorpdfstring{$D^{(n+1,n-1)}$}{D\textasciicircum(n+1,n-1)}}\label{subsubsec:nearmatch}

Let $n\ge1$ and $\mu=(n+1,n-1)$, a $2$-regular partition of $2n$. A $\mu$-tableau $t$ has $n-1$ columns of length two, with entries $a_i$ above $b_i$, and two columns of length one, with entries $d_1,d_2$. Its column stabilizer is the product of the groups $\{1,\tr{a_i}{b_i}\}$, so
\[
e_t\equiv\sum_{R\in\mathcal T(t)}\{R\}\pmod2,
\]
where $\mathcal T(t)$ is the set of $(n+1)$-subsets $R$ that contain $d_1$ and $d_2$ and exactly one of $a_i,b_i$ for each $i$. Put
\[
\theta(t)=\{\{a_i,b_i\}:1\le i\le n-1\}\cup\{\{d_1,d_2\}\}\in\M_n.
\]
The $n-1$ columns of length two of $t$ form a matching of $V\setminus\{d_1,d_2\}$, which we call the \emph{near-matching} of $t$, and $\theta(t)$ completes it by the edge $\{d_1,d_2\}$. Then $\theta(\sigma t)=\sigma\theta(t)$, and $\theta$ is onto, since any edge of a matching can be placed in the two columns of length one.

\begin{lemma}\label{lem:nearmatch}
For $\mu$-tableaux $t,t'$, $\pr{e_t,e_{t'}}\equiv H_{2n}[\theta(t),\theta(t')]\pmod2$.
\end{lemma}

\begin{proof}
Modulo $2$, $\pr{e_t,e_{t'}}$ is the number of sets in $\mathcal T(t)\cap\mathcal T(t')$, since every coefficient is $\pm1$. Let $d_1',d_2'$ be the entries of the two columns of length one of $t'$. We call $d_1,d_2$ the singletons of $t$ and $d_1',d_2'$ the singletons of $t'$. Let $\mathcal G$ be the multigraph on $V$ whose edges are the $n-1$ pairs $\{a_i,b_i\}$ of $t$ and the $n-1$ pairs of $t'$. A vertex lies on at most one pair of each tableau, so its degree is at most $2$, and every component of $\mathcal G$ is a cycle that alternates between pairs of $t$ and of $t'$, or a path. A vertex has degree less than $2$ exactly when it is a singleton of $t$ or of $t'$. A vertex that is a singleton of both tableaux is a path of length $0$. Each path has two ends, counted twice for a path of length $0$, and the ends are exactly the four singleton positions $d_1,d_2,d_1',d_2'$. So $\mathcal G$ has exactly two path components.

Let $R\in\mathcal T(t)\cap\mathcal T(t')$. Then $R$ contains all four singletons and exactly one end of every edge of $\mathcal G$, so along each component the vertices alternate in and out of $R$. On a path both ends are singletons and lie in $R$, so the path has even length. Call a path \emph{through} if it joins a singleton of $t$ to a singleton of $t'$. A path of length $0$ is a common singleton and is through. A path of positive length alternates between pairs of $t$ and of $t'$, and it leaves an end that is a singleton of $t$ along a pair of $t'$, and an end that is a singleton of $t'$ along a pair of $t$. So a path of positive length is through exactly when its length is even. Therefore $\mathcal T(t)\cap\mathcal T(t')$ is empty unless both paths are through, and then $R$ is determined on the paths and has two choices on each cycle. This gives
\[
\pr{e_t,e_{t'}}\equiv\begin{cases}2^{c(\mathcal G)}&\text{if both paths are through},\\0&\text{otherwise},\end{cases}
\]
where $c(\mathcal G)$ is the number of cycles of $\mathcal G$.

On the other side, $\theta(t)\cup\theta(t')$ is $\mathcal G$ with the edges $\{d_1,d_2\}$ and $\{d_1',d_2'\}$ added. If both paths are through, these two edges close the two paths into one cycle. Otherwise one path joins $d_1$ to $d_2$ and the other joins $d_1'$ to $d_2'$, and each closes into its own cycle. So $\kappa(\theta(t)\cup\theta(t'))$ is $c(\mathcal G)+1$ in the first case and $c(\mathcal G)+2$ in the second. It equals $1$ exactly when $\mathcal G$ has no cycle and both paths are through, which is exactly when the right side of the display is odd.
\end{proof}

\begin{theorem}\label{thm:step4}
Let $n\ge1$ and $\mu=(n+1,n-1)$. The rule $e_t\mapsto\varphi(\theta(t))$ extends to a surjective $\F S_{2n}$-homomorphism $\theta_*\colon S^\mu_{\F}\to W_n$ with kernel $S^\mu_{\F}\cap S^{\mu\perp}_{\F}$. Hence
\[
W_n\cong S^\mu_{\F}/(S^\mu_{\F}\cap S^{\mu\perp}_{\F})=D^{(n+1,n-1)},
\]
and $W_n\cong\overline{\mathcal L}^1/\overline{\mathcal L}^2$ is absolutely irreducible.
\end{theorem}

\begin{proof}
Let $x\in S^\mu_{\F}$ and write $x=\sum_{t}c_{t}e_t$ in any way. For $Q\in\M_n$ choose $t'$ with $\theta(t')=Q$. By Lemma~\ref{lem:nearmatch}, the coordinate at $Q$ of $\sum_{t}c_t\varphi(\theta(t))$ is
\[
\sum_{t}c_{t}H_{2n}[\theta(t),Q]=\sum_{t}c_t\pr{e_t,e_{t'}}=\pr{x,e_{t'}} .
\]
The right side depends only on $x$, so $\theta_*(x)=\sum_{t}c_t\varphi(\theta(t))$ is well defined, and it is linear. The same computation with an arbitrary $t''$ in place of $t'$ shows that $\pr{x,e_{t''}}$ is the coordinate of $\theta_*(x)$ at $\theta(t'')$. So $\theta_*(x)=0$ if and only if $\pr{x,e_{t''}}=0$ for every $\mu$-tableau $t''$, that is, $x\in S^{\mu\perp}_{\F}$. Equivariance follows from $\sigma e_t=e_{\sigma t}$, $\theta(\sigma t)=\sigma\theta(t)$ and $\rho_W(\sigma)\varphi(P)=\varphi(\sigma P)$. The image contains every fingerprint because $\theta$ is onto, so $\theta_*$ is onto $W_n$. Since $\mu$ is $2$-regular, the quotient is $D^\mu$ by~\cite[Def.~11.2]{James78}. Absolute irreducibility is Theorem~\ref{thm:burnside}, and also follows from~\cite[Thm.~4.9]{James78}. The last isomorphism is Theorem~\ref{thm:step3}(iv).
\end{proof}

Within this subsection, the rank theorem is used only for the identification with $\overline{\mathcal L}^1/\overline{\mathcal L}^2$ and for dimensions.

\needspace{6\baselineskip}
\begin{corollary}\label{cor:benson}
Let $n\ge1$.
\begin{enumerate}
\item $S^{(n,n)}_{\F}$ has a unique maximal submodule, and its quotient is $D^{(n+1,n-1)}$.
\item $\dim_{\F}D^{(n+1,n-1)}=2^{n-1}$.
\end{enumerate}
\end{corollary}

\begin{proof}
(i) Let $\Psi\colon M^\mu_{\F}\to M^{(n,n)}_{\F}$ send $\{R\}$ to $\sum_{r\in R}\{R\setminus\{r\}\}$. It is $S_{2n}$-equivariant. Apply it to $e_t\equiv\sum_{R\in\mathcal T(t)}\{R\}$. A set $R\setminus\{r\}$ with $r\in\{d_1,d_2\}$ contains exactly one of $d_1,d_2$ and one point of each pair, so it is a transversal of $\theta(t)$, and each such transversal arises from exactly one pair $(R,r)$. A set $R\setminus\{r\}$ with $r$ in a pair $\{a_i,b_i\}$ contains both $d_1,d_2$ and misses both $a_i,b_i$, and it arises from exactly two pairs $(R,r)$. So $\Psi(e_t)\equiv v_{\theta(t)}\pmod2$. Since $\theta$ is onto, $\Psi$ maps $S^\mu_{\F}$ onto $S^{(n,n)}_{\F}$ by Lemma~\ref{lem:vP}. By Theorem~\ref{thm:step4} the quotient $S^\mu_{\F}/(S^\mu_{\F}\cap S^{\mu\perp}_{\F})$ is $W_n\ne0$, so $S^\mu_{\F}\cap S^{\mu\perp}_{\F}$ is the unique maximal submodule of $S^\mu_{\F}$ by~\cite[Thm.~4.9]{James78}. The kernel of $\Psi$ is a proper submodule, so it lies in $S^\mu_{\F}\cap S^{\mu\perp}_{\F}$. Hence $S^{(n,n)}_{\F}\cong S^\mu_{\F}/\ker\Psi$ has the unique maximal submodule $(S^\mu_{\F}\cap S^{\mu\perp}_{\F})/\ker\Psi$, with quotient $D^\mu$. (ii) is Theorem~\ref{thm:step4} with Theorem~\ref{thm:rank}.
\end{proof}

In degree $2n$, part (i) is~\cite[Lemma~5.2]{Benson88} and part (ii) is~\cite[Lemma~5.3]{Benson88}, which Benson derives from James's work on two-part partitions in characteristic $2$. Read in the other direction, the rank theorem $\rank_{\F}H_{2n}=2^{n-1}$ of Cygan, Kratsch and Nederlof is equivalent to $\dim D^{(n+1,n-1)}=2^{n-1}$.

\begin{remark}\label{rem:specht}
For $n=2$, $S^{(2,2)}_{\F}=D^{(3,1)}=W_2$ has dimension $2$ and $\overline{\mathcal L}^2=0$. For $n=1$, $\mu=(2)$ and $W_1=D^{(2)}$ is the trivial module.
\end{remark}
